%% Finite editorial proof supplement; routing and admission are integrator-owned.
\section{Boolean Prime Ideals and Ultrafilter Extension}
\label{sec:p01-boolean-prime-ideal-and-ultrafilters}

The Ultrafilter Lemma of Definition~\ref{def:ultrafilter-lemma} extends a
proper family of subsets to one deciding every subset and its complement.
Its Boolean-algebra formulation is the Boolean Prime Ideal principle.
The equivalence is a theorem of \(\mathsf{ZF}\) about two additional
principles: it does not make either principle a theorem of \(\mathsf{ZF}\),
and it does not identify either with the Axiom of Choice. This treatment gives
the constructions, including the unrestricted propositional compactness
bridge cited in Theorem~\ref{thm:p02-arbitrary-prop-compactness-bpi}.
The algebraic vocabulary follows \cite[Sections 2.2 and 5.4]{Koppelberg1989};
the proofs in the corresponding annex are separately worded ZF arguments.

\subsection{Properness, duality and the two-element algebra}
\label{sec:p01-bpi-algebraic-vocabulary}

\begin{hbdefinition}{def:p01-bpi-boolean-algebra}{Boolean algebra and local notation}
A Boolean algebra \(B\) here is a set with operations
\(\wedge,\vee,\neg\), constants \(0_B,1_B\), and the bounded distributive
lattice and complement laws. Write \(a\leq_B b\) iff
\(a\wedge b=a\). It is nontrivial iff \(0_B\ne1_B\).
The two-element Boolean algebra \(\mathbf2=\{0,1\}\) has its usual
Boolean operations. A homomorphism \(h:B\to\mathbf2\) preserves both
constants, meet, join and complement. These operations and constants are
algebraic, not logical connectives in a theory under consideration.
Koppelberg's \(a\cdot b\), \(a+b\) and \(-a\) denote, respectively,
\(a\wedge b\), \(a\vee b\) and \(\neg a\); they are not numerical
multiplication, addition and subtraction. A finite empty meet is \(1_B\)
and a finite empty join is \(0_B\).
\end{hbdefinition}

\begin{hbdefinition}{def:p01-bpi-filter-ideal}{Filters, ideals, and the two kinds of complement}
A filter \(F\subseteq B\) contains \(1_B\), is upward closed under
\(\leq_B\), and is closed under binary meets. It is proper iff
\(0_B\notin F\). An ideal \(I\subseteq B\) contains \(0_B\), is
downward closed, and is closed under binary joins; it is proper iff
\(1_B\notin I\). A prime ideal is a proper ideal \(P\) such that
\(a\wedge b\in P\) implies \(a\in P\) or \(b\in P\).
An ultrafilter is a maximal proper filter under inclusion.

The dual of a subset \(E\subseteq B\) is
\(d(E)=\{\neg a:a\in E\}\). This is its image under Boolean
complement, not its set-theoretic complement \(B\setminus E\).
For a set \(S\), a proper set filter is a proper Boolean filter in
\(\mathcal P(S)\), with meet \(\cap\), join \(\cup\), complement
\(S\setminus A\), bottom \(\varnothing\) and top \(S\).
A principal set ultrafilter means the point filter
\(U_s=\{A\subseteq S:s\in A\}\) for some \(s\in S\).
Equivalently it is an ultrafilter generated by one subset of \(S\):
the boundary proof below justifies why its generator must be a singleton.
\end{hbdefinition}

\begin{hblemma}{thm:p01-bpi-duality}{proved=annex:p01-bpi-duality}{Choice-free filter and ideal calculus}
For \(E\subseteq B\), the filter generated by \(E\) consists of those
\(b\) above a meet of finitely many members of \(E\); it is proper iff
every such meet is nonzero, including the empty meet. The dual statement
holds for ideals and finite joins. The operation \(d\) is an
inclusion-preserving bijection between filters and ideals, preserving
properness.

For a proper filter \(U\), the following are equivalent: maximality;
membership of exactly one of \(a,\neg a\) for every \(a\in B\);
and primeness in the filter sense
\(a\vee b\in U\Rightarrow a\in U\text{ or }b\in U\).
They are also equivalent to its characteristic function
\(\chi_U:B\to\mathbf2\), defined by \(\chi_U(a)=1\) iff
\(a\in U\) and zero otherwise, being a homomorphism.
For such \(U\), \(d(U)=B\setminus U\) is a prime ideal.
Conversely, every prime ideal is the complement of an ultrafilter and is
maximal among proper ideals. None of these assertions uses Choice.
\end{hblemma}

\begin{hbdefinition}{def:p01-bpi-quotient}{Symmetric difference and the ideal quotient}
Put
\[
 a\mathbin{\triangle}b=(a\wedge\neg b)\vee(\neg a\wedge b),
 \qquad
 a\sim_I b\quad\text{iff}\quad a\mathbin{\triangle}b\in I.
\]
For an ideal \(I\), the intended quotient has underlying set
\(B/I=\{[a]_I:a\in B\}\), where
\([a]_I=\{b\in B:a\sim_I b\}\), and operations induced by those
of \(B\). The next result establishes well-definedness before the quotient
is used; no selection of one representative per class is stipulated.
Koppelberg writes this ideal congruence locally as \(a\equiv b\);
\(\sim_I\) here names that same relation with its ideal parameter
explicit. The quoted source use of \(\equiv\) is a congruence, not the
handbook's elementary-equivalence relation between structures.
\end{hbdefinition}

\begin{hblemma}{thm:p01-bpi-quotient}{proved=annex:p01-bpi-quotient}{Boolean quotient without representative choices}
For every ideal \(I\), \(\sim_I\) is a Boolean congruence and the
displayed quotient is a Boolean algebra. The quotient map
\(q(a)=[a]_I\) is a surjective homomorphism with kernel \(I\).
The quotient is nontrivial exactly when \(I\) is proper. The inverse
image under \(q\) of a prime ideal of \(B/I\) is a prime ideal of
\(B\) containing \(I\).
\end{hblemma}

\subsection{Equivalent extension principles over ZF}
\label{sec:p01-bpi-extension-equivalence}

\begin{hbdefinition}{def:p01-bpi-principles}{The principles being compared}
\(\mathsf{BPI}\) says that every nontrivial set-sized Boolean algebra
has a prime ideal. Boolean ideal extension says that every proper ideal of
every Boolean algebra is contained in a prime ideal. Boolean filter extension
says that every proper Boolean filter is contained in an ultrafilter.
The set Ultrafilter Lemma restricts the last assertion to algebras
\(\mathcal P(S)\), for arbitrary sets \(S\).
These are sentences or schemes of the ambient set theory, not rules of an
object-language deduction calculus. Existence without a prescribed ideal
and extension of a prescribed ideal are distinguished until their
equivalence is proved.
\end{hbdefinition}

\begin{hblemma}{thm:p01-bpi-finite-local}{proved=annex:p01-bpi-finite-local}{Finite local Boolean evaluations}
In a nontrivial Boolean algebra, a subalgebra generated by a finite list
\(b_1,\ldots,b_n\) is finite. Its nonzero cells
\[
 e_\varepsilon=\bigwedge_{k=1}^{n}b_k^{\varepsilon_k},
 \qquad b_k^1=b_k,\quad b_k^0=\neg b_k,
 \qquad \varepsilon\in\{0,1\}^{n},
\]
partition \(1_B\); every member of the subalgebra is a join of cells.
For every finite \(E\subseteq I\), where \(I\) is a proper ideal,
and every finite list containing \(E\), some nonzero cell satisfies
\(e_\varepsilon\leq_B\neg\bigvee E\). On the generated subalgebra,
\(h(c)=1\) iff \(e_\varepsilon\leq_B c\) is a homomorphism to
\(\mathbf2\) and sends each member of \(E\) to zero. The assertion
holds also for empty lists and empty \(E\), without a choice principle.
\end{hblemma}

\begin{hbtheorem}{thm:p01-bpi-equivalence}{proved=annex:p01-bpi-equivalence}{BPI, Boolean extension, and set ultrafilters}
Over \(\mathsf{ZF}\), the following four principles are equivalent:
\(\mathsf{BPI}\), Boolean ideal extension, Boolean filter extension,
and the Ultrafilter Lemma for arbitrary sets. The equivalence includes
prescribed proper filters and ideals; it is not restricted to countable
algebras, well-orderable sets, principal filters or complete Boolean algebras.
\end{hbtheorem}

The two nontrivial constructions serve different purposes. Quotienting by
a prescribed ideal turns bare prime-ideal existence into extension. In the
opposite direction, assignments to all elements of an algebra need only
satisfy each finite collection of Boolean constraints initially. A set
ultrafilter on the assignment space then yields one homomorphism satisfying
all the constraints. The finite local lemma provides the witnesses needed
for properness; it does not choose a witness for every finite collection
simultaneously.

\begin{hbwarning}{def:p01-bpi-empty-degenerate}{Empty sets and degenerate algebras}
There is no proper filter on \(S=\varnothing\), because its top and
bottom coincide. The set extension principle is therefore vacuous there.
Likewise a Boolean algebra with \(0_B=1_B\) has no proper filter or
ideal, so its extension instances are vacuous. Bare prime-ideal existence
must explicitly exclude that algebra: including it would make the
principle false. The algebra \(\mathbf2\) is nontrivial and has the
prime ideal \(\{0\}\) and ultrafilter \(\{1\}\).
\end{hbwarning}

\subsection{The complete propositional compactness bridge}
\label{sec:p01-bpi-compactness-bridge}

\begin{hbdefinition}{def:p01-bpi-propositional-compactness}{Arbitrary-set-variable propositional compactness}
For a set \(V\), use finite propositional formulas built from variables
\(p_v\), constants \(\bot,\top\), and
\(\neg,\land,\lor\). A valuation is a total function
\(\nu:V\to\mathbf2\), extended by the Boolean truth tables.
Compactness for arbitrary variable sets says: whenever
\(\Gamma\) is a set of these formulas and every finite subset of
\(\Gamma\) is satisfied by a valuation on \(V\), one valuation on
\(V\) satisfies all of \(\Gamma\). Use
\(\varphi\Leftrightarrow\psi\) as an abbreviation for
\((\neg\varphi\lor\psi)\land(\neg\psi\lor\varphi)\), whose value
is one exactly when the two formula values agree. This does not add
infinitary formulas. The constants make the case
\(V=\varnothing\) explicit.
\end{hbdefinition}

\begin{hbtheorem}{thm:p01-bpi-propositional-compactness}{proved=annex:p01-bpi-propositional-compactness}{BPI is equivalent to unrestricted propositional compactness}
Over \(\mathsf{ZF}\), the principle in
Definition~\ref{def:p01-bpi-propositional-compactness} is equivalent to
each of the four principles in Theorem~\ref{thm:p01-bpi-equivalence}.
Both directions hold for arbitrary sets of variables and formulas, without
well-ordering either set. The same compactness statement with a fixed
enumerable variable set has a separate choice-free proof and must not be
used as though it were the arbitrary-set principle.
\end{hbtheorem}

The forward construction takes the truth sets of premises inside the set
of all valuations and extends their generated filter. For the converse,
one variable per Boolean element expresses the Boolean operation laws and
the exclusion of a prescribed ideal. This accounts for the choice strength:
a finite formula always has a finite truth table, but there may be no
enumeration supporting a successive decision of every variable in an
arbitrary set. The fixed binary-tree route in
\cref{sec:p06-cardinals-choice} is not being strengthened here.
Nor do finite truth tables supply the missing valuation: the decision
procedures for propositional logic act on finite encoded inputs over
coded atoms, and \cref{sec:p02-w02-p02-017-finite-input-scope} shows that
they neither decide infinite premise sets nor construct the valuation
that compactness asserts.
Adjoining truth constants does not change compactness for nonempty alphabets:
one fixed variable supplies a tautology and contradiction in their place.
For the empty alphabet the constants-only truth evaluation is unique.

\subsection{Choice strength and concrete boundaries}
\label{sec:p01-bpi-choice-strength}

\begin{hbtheorem}{thm:p01-bpi-ac-implication}{proved=annex:p01-bpi-ac-implication}{Choice supplies ultrafilter extension}
In \(\mathsf{ZF}+\mathsf{AC}\), every proper Boolean filter extends
to an ultrafilter, so \(\mathsf{BPI}\) holds. More precisely, Zorn's
lemma suffices for this extension. The sole maximality principle used in
this proof is supplied by Choice through
Theorem~\ref{thm:p06-ac-to-zorn}; no such principle is used in the ZF
equivalence proofs above.
\end{hbtheorem}

\begin{hbtheorem}{thm:p01-bpi-choice-separation}{cited=HalpernLevy1971}{Relative non-implication of Choice}
The historical separation theorem has the relative-consistency form
\[
 \operatorname{Con}(\mathsf{ZF})\ \Longrightarrow
 \operatorname{Con}(\mathsf{ZF}+\mathsf{BPI}+\neg\mathsf{AC}).
\]
Consequently, if \(\mathsf{ZF}\) is consistent, the implication from
\(\mathsf{BPI}\) to \(\mathsf{AC}\) is not provable over
\(\mathsf{ZF}\). This is a cited independence theorem, not a model
construction proved by the elementary filter arguments in this treatment.
\end{hbtheorem}

For a later, distinct primary presentation, Repick\'y starts with a ZFC
ground universe and an appropriate generic extension, and proves BPI in a
symmetric inner model where Choice fails \cite{Repicky2015BPI}.
The existence and ZF properties of that model and its failure of Choice are
imported there from the cited set-theoretic construction; his finite
continuity argument is not a choice-free construction of a prime ideal in
every universe. Nor does a theorem saying that AC is sufficient prove AC
necessary. No comparison with Dependent Choice or Countable Choice is
deduced from this treatment's equivalences.

\begin{hbcorollary}{thm:p01-bpi-set-boundaries}{proved=annex:p01-bpi-set-boundaries}{Finite principal and infinite cofinite boundaries}
In \(\mathsf{ZF}\), every proper filter on a finite set extends to a
principal ultrafilter, and every ultrafilter on a finite set is principal.
For any infinite set \(S\), its cofinite subsets form a proper filter.
Assuming \(\mathsf{BPI}\), this filter has an ultrafilter extension
containing no finite subset of \(S\), hence no singleton and therefore
not principal. No countably infinite subset of \(S\) is required.
\end{hbcorollary}

\begin{hbexample}{def:p01-bpi-worked-finite}{A prescribed filter has several possible extensions}
Let \(S=\{a,b,c\}\) and let \(F\) be the filter generated by
\(\{a,b\}\); thus its members are \(\{a,b\}\) and \(S\).
The two ultrafilter extensions are
\(U_a=\{A\subseteq S:a\in A\}\) and
\(U_b=\{A\subseteq S:b\in A\}\).
Under ideal duality, \(d(F)=\{\varnothing,\{c\}\}\) extends to
\(\mathcal P(S\setminus\{a\})\) or
\(\mathcal P(S\setminus\{b\})\), respectively. These are prime
ideals, the zero kernels of evaluation at the selected point. The filter
does not select a unique extension, and no uniqueness is promised by BPI.
On \(\mathbb N\), by contrast, extending the cofinite filter excludes
every point evaluation: each singleton's complement already belongs to the
filter. That latter existence assertion uses the additional principle.
\end{hbexample}
